\chapter{Extensions nulles et clôture matérielle}
\label{chap:realizaciones-nula-material}
\index{photon!interface physique}\index{électron!clôture matérielle}
\index{extension nulle}\index{cône causal scindé}

\section{\(4\) réalisations de l'état matériel commun : nulle, massive, composée et topologique}
\label{sec:cuatro-modos-realizacion}

Un chiffre isolé détermine une coordonnée, mais n'identifie ni la structure
qui le produit, ni la mémoire qu'il conserve, ni l'opération qui le transforme en
grandeur. HMT se situe avant cette projection : son objet primaire est une
extension compatible qui conserve les lois arithmétiques, l'orientation
TRIT, les quotients, la phase et la mémoire de frontière. La mécanique dimensionnelle
spécifie ensuite quelle composante se réalise comme grandeur, laquelle demeure
comme mémoire et quelle transformation produit l'observable.

Les développements de cette partie distinguent quatre modes par leur opération de
fermeture :
\begin{enumerate}
 \item une extension ouverte et nulle transporte une perturbation entre
 frontières ;
 \item une extension fermée et persistante acquiert une norme intérieure et définit
 le candidat matériel ;
 \item une famille d'extensions pondérée par l'action définit un ensemble
 thermique ;
 \item une connexion qui compare des sections en des points distincts fournit le
 codomaine de la réalisation gravitationnelle.
\end{enumerate}
La source commune n'identifie pas les quatre codomaines. Les applications physiques
correspondantes restent discontinues :
\[
 \mathfrak R_\gamma:
 \Omega_{\HMT}^{\rm surv}\dashrightarrow\mathcal H_{\rm rad},
 \qquad
 \mathfrak R_m:
 \Omega_{\HMT}^{\rm surv}\dashrightarrow\mathcal H_{\rm mat},
\]
\[
 \mathfrak R_\Theta:
 \mathcal P(\Omega_{\HMT}^{\rm surv})\dashrightarrow\mathcal S_{\rm th},
 \qquad
 \mathfrak R_{\rm grav}:
 \Omega_{\HMT}^{\rm surv}\dashrightarrow\mathsf{Cartan}_{1,3}.
\]
Leur unité réside dans la généalogie ; leur différence, dans la lecture et dans les
structures supplémentaires exigées par chaque réalisation.

\section{Propagation nulle entre frontières}
\label{sec:propagacion-nula-fotonica}

Soit \(B\) une base de transport. Une \emph{section nulle} de signe
\(\pm\) est une application
\[
 s_\pm:B\longrightarrow\mathcal A_{\rm sp}P_\pm
\]
telle que \(s_\pm(x)\in\mathcal A_{\rm sp}P_\pm\) pour tout \(x\in B\).
Par la \cref{prop:descomposicion-espectral-partida},
\[
 N_{\rm sp}(s_\pm(x))=0
 \qquad(x\in B).
\]
Si \(\gamma\) est une voie orientée entre deux frontières,
\[
 \partial\gamma=x_{\rm abs}-x_{\rm em},
\]
l'interface photonique proposée attribue à \(\gamma\) une section nulle
étendue. L'adjectif \emph{étendue} indique que son identité appartient au
transport complet et non à un point isolé.

La voie conjuguée s'obtient par inversion,
\[
 \gamma^\vee=C_{\rm rev}\gamma,
 \qquad
 \partial\gamma^\vee=x_{\rm em}-x_{\rm abs}.
\]
Relativement au lecteur d'orientation déclaré, la paire
\((\gamma,\gamma^\vee)\) représente les lectures avancée et retardée. Les
deux branches sont ordonnées séquentiellement ou logées dans des fibres distinctes. Elles ne
s'ajoutent pas automatiquement dans une même fibre : si les deux composantes conjuguées
sont activées simultanément, l'identité exacte
\[
 \boxed{N_{\rm sp}(uP_++vP_-)=uv}
\]
montre que, pour \(u,v>0\), la section quitte la frontière nulle et entre dans
l'intérieur positif.

\subsection{Longueur, aller et retour}

La carte dimensionnelle de vitesse,
\(\widehat c=e^{-\mathcal E}\), et une base positive \(u_L\) de la droite de
longueur étant fixées, on définit
\[
 L_{\rm int}(\gamma)=\widehat c\,|\gamma|\,u_L.
\]
L'inversion conserve le nombre de pas :
\[
 L_{\rm int}(\gamma^\vee)=L_{\rm int}(\gamma),
 \qquad
 L_{\leftrightarrow}=2L_{\rm int}(\gamma).
\]
Ces identités sont cinématiques. Les transformer en propagation
électromagnétique exige qu'une même application reproduise la loi de Gauss, la
dynamique de jauge, la dispersion nulle et la réduction à deux degrés
transversaux.

\subsection{Entropie d'intrication entre les \(2\) lectures}
\label{sec:entrelazamiento-ida-retorno}

L'inversion d'une voie ne suffit pas pour affirmer une intrication quantique. Une
formulation minimale exige d'abord une complétion de Hilbert. Soient
\(\mathcal H_{\rm adv}\) et \(\mathcal H_{\rm ret}\) deux espaces complexes
avec des familles orthonormales
\(\{|i\rangle_{\rm adv}\}_{i\in I}\) et
\(\{|i\rangle_{\rm ret}\}_{i\in I}\), où \(I\) est fini. Pour une
probabilité \(p=(p_i)_{i\in I}\), nous définissons
\[
 |\Psi_p\rangle
 :=\sum_{i\in I}\sqrt{p_i}\,
 |i\rangle_{\rm adv}\otimes|i\rangle_{\rm ret}.
 \tag{E1}\label{eq:estado-entrelazado-ida-retorno}
\]

\begin{proposition}[Entropie bicouche de l'aller et du retour]
\label{prop:entropia-entrelazamiento-ida-retorno}
Les matrices réduites de l'état
\(\rho_p=|\Psi_p\rangle\langle\Psi_p|\) sont
\[
 \rho_{\rm adv}
 =\sum_{i\in I}p_i|i\rangle_{\rm adv}\langle i|,
 \qquad
 \rho_{\rm ret}
 =\sum_{i\in I}p_i|i\rangle_{\rm ret}\langle i|,
 \tag{E2}\label{eq:reducidas-ida-retorno}
\]
et possèdent la même entropie de von Neumann :
\[
 \boxed{
 S_{\rm ent}(p)
 =-\operatorname{Tr}(\rho_{\rm adv}\log\rho_{\rm adv})
 =-\sum_{i\in I}p_i\log p_i.}
 \tag{E3}\label{eq:entropia-ida-retorno}
\]
On adopte ici la convention \(0\log0:=0\). L'état est séparable
exactement lorsqu'une seule probabilité vaut un.
\end{proposition}

\begin{proof}
En effectuant la trace partielle, l'orthogonalité annule les termes avec
\(i\ne j\). Les valeurs propres non nulles de chaque matrice réduite sont donc
les \(p_i\), ce qui donne \eqref{eq:entropia-ida-retorno}. Le rang de
Schmidt est un exactement dans le cas déterministe.
\end{proof}

La proposition formalise, une fois la complétion choisie, une entropie entre
les lectures avancée et retardée. Elle ne déduit pas les probabilités \(p_i\),
n'identifie pas ces lectures aux propagateurs avancé et retardé d'une
théorie des champs et ne transforme pas une corrélation classique de voies en
intrication. La construction d'un isomorphisme qui envoie la mémoire HMT
sur des états physiques bipartites, préserve l'évolution et fixe \(p\) appartient à
l'\textsc{interface physique}. L'entropie nulle d'une distribution
déterministe constitue le contrôle négatif immédiat.

\begin{statusbox}
La proposition est
\textsc{démontrée avec une structure de départ explicite} : elle présuppose la
complétion de Hilbert, les bases orthonormales et la distribution \(p\). La construction
de l'opérateur d'entrelacement qui réalise physiquement les deux couches est une
\textsc{interface physique} ; elle ne fait pas partie de la proposition.
\end{statusbox}

\begin{definition}[Interface photonique de HMT--MD]
\label{def:interfaz-fotonica-hmt-md}
Une réalisation photonique de HMT--MD est une section ouverte portée par
une direction nulle de l'algèbre scindée, transportée entre des frontières
conjuguées et protégée par une symétrie qui conserve exactement la norme
nulle.
\end{definition}

Les idempotents \(P_\pm\) ne sont pas, à eux seuls, des photons physiques ni les deux
polarisations de Maxwell. La réalisation complète exige un
opérateur d'entrelacement
\[
 \mathfrak R_\gamma:
 \partial_{\rm null}\mathcal C_{\rm sp}^{+}
 \dashrightarrow\mathcal H_{\rm Maxwell}^{\rm phys}
\]
qui préserve le produit scalaire, l'évolution et l'action de jauge, et qui
sélectionne exactement deux polarisations transversales en \(3+1\)
dimensions.

\begin{criterion}[Secteur photonique]
\label{crit:sector-fotonico}
L'identification physique est réfutée si l'évolution quitte
\(\partial_{\rm null}\mathcal C_{\rm sp}^{+}\), engendre une masse sans le mécanisme
de brisure déclaré, ne reproduit pas la dispersion nulle ou ne laisse pas exactement
deux polarisations physiques après la réduction de jauge.
\end{criterion}

\section{Clôture matérielle et centre électronique}
\label{sec:clausura-material-electron}

Le couplage simultané de deux directions nulles conjuguées produit, pour
\(u,v>0\),
\[
 N_{\rm sp}(uP_++vP_-)=uv>0.
\]
Aucune composante isolée n'a de norme non nulle ; la norme intérieure
naît de leur couplage. Cet énoncé algébrique n'additionne pas des masses
préexistantes et n'identifie pas encore les directions nulles à des photons
\emph{sur couche de masse}.

   La classification structurelle de HMT contient un unique point fixe de l'inversion cellulaire,
\[
 (5,5).
\]
L'origine relative d'action de l'électron est \(v_e=0\). Dans la réalisation
scindée normalisée, on fixe
\[
 Z_e=P_++P_-=I,
 \qquad
 N_{\rm sp}(Z_e)=1.
\]
Le point \((5,5)\), l'origine \(v_e=0\), la signature brute
\((24,0,12,0,-12)\) et le mot transversal \(555555\) sont quatre objets
distincts.

\begin{definition}[Candidat de clôture électronique]
\label{def:candidato-clausura-electronica}
La clôture résonante de modes nuls électromagnétiques conjugués se
formalise comme le programme qui cherche une section matérielle centrale normalisée
obtenue par le couplage de deux modes nuls conjugués, stabilisée par
une clôture de voie et par son holonomie.
\end{definition}

L'existence de l'extension, du retour et du centre de l'atlas ne prouve pas
la compacité minimale ni la stabilité de la fermeture. Le théorème dynamique
restant à établir doit construire
\[
 \gamma_e=\xi_1\circ\cdots\circ\xi_N,
 \qquad
 \partial\gamma_e=0,
 \qquad
 \varepsilon(\gamma_e)=-1,
\]
où chaque \(\xi_j\) est localement nul, le bord total s'annule et
\(\varepsilon\) est l'holonomie de signe. La masse serait associée au cocycle
global de clôture, non à une somme de masses nulles locales. Ce programme ne
rouvre ni l'électron central \((5,5)\), ni l'atlas, ni la chaîne
extension--voie--registre--caractère déjà construits.

\begin{proposition}[Critère relativiste d'intérieur causal]
\label{prop:cierre-nulos-interior-causal}
Soient \(p_1,\ldots,p_N\) des impulsions nulles futures dans un espace de Minkowski de
signature temporelle positive, \(p_i^2=0\), et soit
\[
 P=\sum_{i=1}^{N}p_i.
\]
Si tous les \(p_i\) sont colinéaires, alors \(P^2=0\). Si au moins deux sont
futurs et non colinéaires, alors \(P^2>0\), et l'on peut définir
\[
 m^2c^2=P^2.
\]
\end{proposition}

\begin{proof}
Pour des vecteurs nuls futurs, on a
\(p_i\cdot p_j\geq0\), avec égalité exactement lorsqu'ils sont colinéaires.
Comme
\[
 P^2=\sum_i p_i^2+2\sum_{i<j}p_i\cdot p_j,
\]
la première somme est nulle. Tous les produits croisés s'annulent dans le cas
colinéaire ; s'il existe une paire non colinéaire, au moins l'un est strictement
positif.
\end{proof}

Ainsi, une voie fermée n'est pas massive par la seule topologie, et une voie
ouverte n'est pas nulle du seul fait qu'elle possède un bord. La réalisation exige que
la fermeture géométrique et la somme des impulsions entrent conjointement dans
l'intérieur causal.

\subsection{Möbius, conjugaison et occupation ultérieure}

Sous les hypothèses de relèvement et d'holonomie du
\cref{thm:dos-cuatro-pi-estadistica}, un tour inverse l'orientation et
deux la restaurent. Le ruban de Möbius organise la particule et l'antiparticule comme des
orientations conjuguées d'un même type de fermeture. Leur identification physique
exige une représentation qui inverse la charge, préserve la norme et
entrelace l'évolution. La géométrie de Möbius n'implique pas à elle seule
l'algèbre d'occupation fermionique ; ses relations opératoires sont
introduites comme hypothèses opératoires dans le chapitre suivant.

\begin{criterion}[Fermeture électronique]
\label{crit:cierre-electronico}
La réalisation est réfutée si le centre \((5,5)\) ne se relève pas en une
extension fermée, si l'holonomie ne sélectionne pas le secteur impair, si
l'hamiltonien ne possède pas de solution liée stable et unique à symétrie près, ou
si la charge, la dispersion et le facteur de forme ne concordent pas avec la réalisation
électronique.
\end{criterion}
